\chapter{High-Dimensional Computational Geometry}
\label{ch:high_dimensional}

\section{Geometry beyond Three Dimensions}
\label{sec:highdim_intro}

Almost every algorithm in this book was presented for the plane or for three
dimensions, and its analysis relied on a geometric fact that quietly fails in
high dimensions: a straight-line planar subdivision has only $O(n)$ edges, a
triangulation of a point set in the plane has only $O(n)$ triangles, and two
lines in the plane either meet or are parallel. When the dimension $d$ is no
longer a fixed constant but a quantity comparable to $n$, or even larger, the
combinatorial explosion of the objects that computational geometry studies
changes the rules of the game. This chapter studies what happens to the
fundamental structures of the book --- convex hulls, Delaunay triangulations,
nearest-neighbor search, and range searching --- as the dimension grows, and it
develops the two families of ideas that make high-dimensional computation
tractable in practice: geometric reduction, led by the
Johnson--Lindenstrauss lemma, and data-dependent hashing, led by
locality-sensitive hashing. Along the way we meet the
\index{curse of dimensionality}curse of dimensionality, the concentration of
measure that concentrates volume and distance where intuition does not expect
them, and the reasons why approximate answers, not exact ones, are the right
target in high dimension.

\begin{definition}[Dimension regime]
\label{def:highdim_regime}
The dimension $d$ is called
\index{fixed dimension}\emph{fixed} when it is a constant independent of the
input size $n$, so that constants of the form $O(d^{O(d)})$ or $O(2^d)$ may
hide inside the big-$O$ of the running time. It is called \emph{growing} or
\emph{high} when $d$ is an input parameter that may be comparable to $n$; the
extreme case $d \approx n$ appears, for example, in the linear programming
setting of Chapter~\ref{ch:linear_programming} and in statistical
applications.
\end{definition}

The contrast with the rest of this book is sharpest when we count cells. In
the plane the Voronoi diagram of $n$ points has $O(n)$ vertices, and the
combinatorial complexity of the convex hull is $O(n)$ (Chapters
\ref{ch:proximity} and \ref{ch:convex_hulls}). In $\mathbb{R}^d$ the analogous
counts grow like $n^{\lfloor d/2 \rfloor}$, which is linear only for $d \le 3$.
The complexity barrier is the subject of Section~\ref{sec:highdim_hulls} and
Section~\ref{sec:highdim_voronoi_delaunay}.

\begin{theorem}[Complexity of convex hulls]
\label{thm:highdim_hull_complexity}
The number of facets of the convex hull of $n$ points in $\mathbb{R}^d$ is at
most $O(n^{\lfloor d/2 \rfloor})$, and this bound is tight: for every fixed
$d$ there are point sets whose hull has
$\Theta(n^{\lfloor d/2 \rfloor})$ facets. The same bound holds for the number
of $d$-simplices of the Delaunay triangulation and for the number of cells of
the Voronoi diagram.
\end{theorem}

Section~\ref{sec:highdim_concentration} explains the counterintuitive
geometry of high-dimensional balls, Gaussians, and convex bodies, and
Section~\ref{sec:highdim_curse} converts that geometry into an algorithm
analysis: why nearest neighbors lose their meaning, why range trees and
$k$d-trees degrade to linear scans, and why worst-case approximation is hard.
The two positive threads run through Sections
\ref{sec:highdim_dimensionality_reduction}, \ref{sec:highdim_jl}, and
\ref{sec:highdim_random_projections} (random embeddings that compress the data
while approximately preserving pairwise distances) and Sections
\ref{sec:highdim_nearest_neighbor} and \ref{sec:highdim_curse} (hashing schemes
whose query cost is a sublinear power of $n$). Section~\ref{sec:highdim_ml}
ties the techniques to the machine learning pipeline: $k$-nearest-neighbor
classifiers, feature maps, and the algorithmic geometry of learning.

Throughout the chapter the norm $\|\cdot\|$ is the Euclidean norm unless stated
otherwise, and a point set is written $P \subset \mathbb{R}^d$ with
$|P| = n$.

\index{high-dimensional computational geometry}

\section{High-Dimensional Convex Hulls}
\label{sec:highdim_hulls}

\subsection{1. Overview}

The convex hull problem is the most basic one in computational geometry, and
its high-dimensional version is the cleanest illustration of why dimension is
a cost parameter. Chapter~\ref{ch:convex_hulls} showed that the planar hull of
$n$ points can be computed in $O(n \log n)$ time, matching the information
theoretic lower bound, and that the hull has at most $2n$ vertices. In
$\mathbb{R}^d$ the hull can have $\Theta(n^{\lfloor d/2 \rfloor})$ facets
(Theorem~\ref{thm:highdim_hull_complexity}); for $d = 4$ this is quadratic in
$n$, and for $d$ around $\log n$ it is $n^{\Omega(\log n / \log\log n)}$,
larger than any polynomial. Any algorithm that writes the hull explicitly must
therefore spend at least $n^{\lfloor d/2 \rfloor}$ time, so the only
meaningful question is whether the combinatorial explosion of the output is
the \emph{only} cost. The answer is yes: randomized incremental construction
achieves expected time proportional to the output size, plus an $O(n \log n)$
sorting term, for every fixed $d$.

\begin{theorem}[Randomized incremental hulls]
\label{thm:highdim_hull_randomized}
For any fixed dimension $d$, the convex hull of $n$ points in general position
in $\mathbb{R}^d$ can be computed in expected
$O(n \log n + n^{\lfloor d/2 \rfloor})$ time. For $d \le 3$ the bound becomes
$O(n \log n)$.
\end{theorem}

The proof of Theorem~\ref{thm:highdim_hull_randomized} is a high-dimensional
version of the randomized incremental argument of
Chapter~\ref{ch:randomized_algorithms}: insert the points one by one in random
order, and amortize the cost of a point against the facets it destroys. The
analysis depends on a critical fact called the \index{Clarkson--Shor
technique}Clarkson--Shor technique: the expected number of facets destroyed by
a random insertion equals, up to constant factors, the expected number of
facets of a random sample, and the latter can be bounded directly.

\subsection{2. Definitions}

\begin{definition}[Facet, ridge, horizon]
\label{def:highdim_facet}
A \emph{facet} of a $d$-dimensional convex polytope is a
$(d-1)$-dimensional face. A \emph{ridge} is a $(d-2)$-dimensional face. In a
simplicial polytope every facet is a simplex spanned by $d$ vertices and every
ridge is spanned by $d-1$ vertices. Given a point $p$ outside a polytope
$H$, a facet $f$ is \emph{visible from} $p$ if $p$ lies in the open
halfspace outside the supporting hyperplane of $f$. The \emph{horizon} of $p$
on $H$ is the set of ridges that belong to exactly one visible facet.
\end{definition}

In the plane the horizon of an outside point is a pair of vertices and the
visible facets form a contiguous chain along the polygon boundary. In
$\mathbb{R}^d$ the visible facets still form a topological disk on the hull,
bounded by the horizon ridges; this is the only structural fact the algorithm
needs.

\begin{definition}[Cyclic polytope]
\label{def:highdim_cyclic}
The \index{cyclic polytope}\emph{cyclic polytope} $C_d(n)$ is the convex hull
of $n$ points on the moment curve
$\{ (t, t^2, \dots, t^d) : t \in \mathbb{R} \}$. It is \emph{neighborly}:
every $\lfloor d/2 \rfloor$ vertices form a face. Its number of facets is
$O(n^{\lfloor d/2 \rfloor})$, and it is the extremal polytope achieving the
maximum number of facets among all $d$-polytopes with $n$ vertices.
\end{definition}

Cyclic polytopes show that the upper bound of
Theorem~\ref{thm:highdim_hull_complexity} is attainable and that the answer to
any algorithmic question about hulls in dimension $d \ge 4$ must be
input-sensitive.

\subsection{3. Theory}

The high-dimensional hull inherits the polarity duality of
Chapter~\ref{ch:convex_hulls}: facets of the hull of $P$ correspond to
vertices of the polar polytope. The structural result underlying all exact
algorithms is the upper bound theorem
(Theorem~\ref{thm:highdim_hull_complexity}), and the algorithmic result is
that the hull can be found within the output bound plus $O(n \log n)$.

A word on lower bounds. Computing the hull of $n$ points in $\mathbb{R}^d$
requires $\Omega(n \log n)$ time (sorting reduces to it) and
$\Omega(n^{\lfloor d/2 \rfloor})$ time (the output may be that large). For
fixed $d$ the randomized incremental algorithm matches both bounds in
expectation; the deterministic algorithm of Chazelle achieves the same
asymptotics and establishes that the problem has optimal worst-case
complexity in every fixed dimension.

\subsection{4. Pseudocode}

Algorithm~\ref{alg:highdim_hull} is the randomized incremental algorithm in
its full generality. It reduces to the planar incremental algorithm of
Chapter~\ref{ch:convex_hulls} for $d = 2$.

\begin{algorithm}[htbp]
\caption{Randomized Incremental Convex Hull in $\mathbb{R}^d$}
\label{alg:highdim_hull}
\begin{algorithmic}[1]
\Require A set $P$ of $n$ points in general position in $\mathbb{R}^d$.
\Ensure The facets of the convex hull of $P$.
\Function{IncrementalHull}{$P$}
  \State shuffle the points of $P$ uniformly at random
  \State let $H$ be the hull of the first $d+1$ points
  \For{each remaining point $p$}
    \State mark as visible every facet of $H$ that has $p$ on its outer side
    \If{no facet of $H$ is visible from $p$}
      \State \textbf{continue} \Comment{$p$ lies inside $H$}
    \EndIf
    \State collect the horizon: ridges of $H$ shared by one visible and one hidden facet
    \For{each horizon ridge $r$}
      \State add the new facet spanned by $r$ and $p$
    \EndFor
    \State delete all visible facets from $H$
  \EndFor
  \State \Return $H$
\EndFunction
\end{algorithmic}
\end{algorithm}

The inner loop over facets is the whole algorithm in the plane; the only
genuinely $d$-dimensional bookkeeping is locating the visible facets and
walking their boundary to extract the horizon. Both are handled by a facet
incidence graph maintained incrementally. Each inserted point either lies
inside (and is discarded) or destroys the visible facets and creates the cone
of new facets over the horizon.

\subsection{5. Parameter Selection}

The algorithm has no numerical parameters, but its analysis depends on the
assumption of general position: no $d+1$ points coplanar with a facet and no
point lying on a supporting hyperplane. In practice one either perturbs the
input symbolically or, better, works with an exact arithmetic layer that
evaluates the orientation predicate of the facet with respect to $p$; this is
the $d$-dimensional determinant
$\operatorname{orient}(p_0,\dots,p_d)$, the natural generalization of the
predicate used throughout this book. Exact computation is essential because a
single misclassified visibility test invalidates the horizon walk and can
destroy the combinatorial structure.

\subsection{6. Complexity}

Each point inserted destroys a set of facets and creates new ones; over the
whole run the number of created facets is bounded by the expected output size.
The Clarkson--Shor argument shows that the expected number of facets created
by the $i$-th insertion is $O(i^{\lfloor d/2 \rfloor - 1})$, so the total
expected number of created facets is $O(n^{\lfloor d/2 \rfloor})$, and each
creation is charged a constant amount of work. The expected running time is
therefore
\[
O(n \log n + n^{\lfloor d/2 \rfloor}),
\]
with the logarithmic term coming only from the initial shuffle and from
accessing each point once. The space is proportional to the number of facets
currently alive, $O(n^{\lfloor d/2 \rfloor})$ in the worst case.

\subsection{7. Usages}

The incremental hull is the engine behind fixed-dimensional Delaunay
triangulations via lifting (Section~\ref{sec:highdim_voronoi_delaunay}),
behind the computation of regular triangulations and secondary polytopes, and
behind the fixed-dimensional linear programming algorithms studied in
Chapter~\ref{ch:linear_programming}, whose randomized incremental structure is
identical. It is also the standard route to the computation of the volume of
convex polytopes in fixed dimension, where one triangulates the hull into
$d$-simplices. For growing dimension the algorithm is unusable, since the
output bound $n^{\lfloor d/2 \rfloor}$ is exponential in $d$; the remedies are
linear programming duality when only a single linear functional is sought
(Section~\ref{sec:highdim_concentration}) and approximation, which is the
subject of the rest of this chapter.

\subsection{8. Testing Methods and Edge Cases}

Edge cases to test: points all inside a small cloud (interior insertions must
be rejected by the visibility test without modifying the hull); points lying on
facet supporting hyperplanes (requires symbolic perturbation); duplicate and
coplanar points (handled by the general position convention); and the
degenerate starting configuration when the first $d+1$ points are affinely
dependent. A useful invariant-based test verifies, after every insertion, that
every facet orientation is consistently outward and that the facet graph is
manifold: every ridge bounds exactly two facets. For validation against
brute force, enumerate all $d+1$-subsets and check that the facets computed are
exactly the subsets whose supporting hyperplane leaves every other point on
one side; this is exponential in $n$ and only feasible for $n \le 20$. The
extremal cases are the cyclic polytopes of
Definition~\ref{def:highdim_cyclic}, which exercise the upper-bound-size
output and the horizon machinery maximally.

\subsection{9. References}

The incremental hull and its Clarkson--Shor analysis appear in
\cite{mulmuley1994}; the deterministic optimal algorithm is due to
\cite{chazelle1993}. The upper bound theorem and cyclic polytopes are treated
in \cite{matousek2002}. The standard planar treatment appears in
\cite{berg2008}, and the fixed-dimensional hull is surveyed in
\cite{deberg2008}.

\section{Voronoi and Delaunay Structures}
\label{sec:highdim_voronoi_delaunay}

\subsection{1. Overview}

Chapter~\ref{ch:proximity} developed the planar Voronoi diagram and its dual,
the Delaunay triangulation. Both objects remain well defined in
$\mathbb{R}^d$: the Voronoi cell of a site is the set of points at least as
close to that site as to any other, and the Delaunay complex connects two
sites when there is a sphere through them with no other site inside. The
combinatorial complexity, however, explodes from $O(n)$ to
$n^{\lceil d/2 \rceil}$, and --- more dramatically --- the Delaunay complex is
no longer necessarily a triangulation of the space. Beyond the plane the
Delaunay complex may have missing cells, and its topology is not in general a
 valid subdivision of $\mathbb{R}^d$. Understanding this failure is essential
 for meshing and for the geometry processing algorithms of
 Chapter~\ref{ch:mesh_analysis} scaled to higher dimensions.

\subsection{2. Definitions}

\begin{definition}[Voronoi cell, Delaunay triangulation]
\label{def:highdim_voronoi}
For a finite site set $P \subset \mathbb{R}^d$ and a site $p \in P$ the
\emph{Voronoi cell} of $p$ is
\[
\operatorname{vor}(p) = \{ x \in \mathbb{R}^d : \|x - p\| \le \|x - q\| \text{ for all } q \in P \}.
\]
The \emph{Delaunay complex} of $P$ is the collection of all subsets
$S \subseteq P$ for which there exists a sphere passing through the points of
$S$ and containing no point of $P$ in its interior. When every Delaunay cell
is a $d$-simplex, the complex is a triangulation of
$\operatorname{conv}(P)$ and is called the \emph{Delaunay triangulation}.
\end{definition}

The $k$-face $(k+1)$-dimensional Voronoi cells are in bijection with the
$(d-k)$-dimensional Delaunay faces, exactly as in the plane, and the whole
structure is captured by the lifting to the paraboloid.

\begin{definition}[Lifting map]
\label{def:highdim_lifting}
The \index{lifting map}\emph{lifting map} sends $x \in \mathbb{R}^d$ to
$\widehat{x} = (x, \|x\|^2) \in \mathbb{R}^{d+1}$, the point on the paraboloid
$y = \|x\|^2$. A sphere with center $c$ and radius $r$ in $\mathbb{R}^d$
lifts to a hyperplane in $\mathbb{R}^{d+1}$ that passes through the lifted
images of all points of the sphere.
\end{definition}

\subsection{3. Theory}

The lifting map is the master reduction. The Delaunay complex of $P$ in
$\mathbb{R}^d$ is the projection of the lower hull of the lifted points
$\{\widehat{p} : p \in P\}$ in $\mathbb{R}^{d+1}$: a $k$-simplex of the
Delaunay complex corresponds to a $(k+1)$-facet of the lower hull. This is the
statement proved in the plane in Chapter~\ref{ch:proximity}, and it extends
verbatim to all dimensions. The consequence is an immediate transfer of every
hull algorithm from Section~\ref{sec:highdim_hulls}:

\begin{theorem}[Complexity of Delaunay and Voronoi]
\label{thm:highdim_delaunay_complexity}
For any fixed dimension $d$, the Delaunay complex of $n$ points in
$\mathbb{R}^d$ has at most $O(n^{\lceil d/2 \rceil})$ $d$-simplices and can be
computed in expected $O(n \log n + n^{\lceil d/2 \rceil})$ time. The Voronoi
diagram has the same complexity up to constant factors.
\end{theorem}

Two structural warnings follow from the same theorem. First, for $d \ge 3$ the
number of $d$-simplices can be superlinear, so the total work is dominated by
the size of the output, which in turn can be made large by lifting cyclic
polytopes. Second, the Delaunay complex need not be a triangulation of
$\operatorname{conv}(P)$: some $d$-cells may be missing, and the complex may
fail to be a manifold. The Delaunay triangulation in $\mathbb{R}^3$ (for
example) exists only for point sets that are \index{inscribed
polytope}inscribable in the right sense; this is a geometric obstruction with
 no planar analogue and is the reason meshing in higher dimension is
 qualitatively harder than the surface meshing of Chapter~\ref{ch:mesh_generation}.

\subsection{4. Pseudocode}

Algorithm~\ref{alg:highdim_delaunay} reduces the Delaunay computation to the
incremental hull of Algorithm~\ref{alg:highdim_hull} applied to lifted points.
It is the same strategy used in the plane, now without the help of the planar
structure.

\begin{algorithm}[htbp]
\caption{Delaunay Complex in $\mathbb{R}^d$ by Lifting}
\label{alg:highdim_delaunay}
\begin{algorithmic}[1]
\Require A set $P$ of $n$ points in general position in $\mathbb{R}^d$.
\Ensure The Delaunay complex of $P$.
\Function{Delaunay}{$P$}
  \State compute $\widehat{P} = \{ (p, \|p\|^2) : p \in P \} \subset \mathbb{R}^{d+1}$
  \State compute the convex hull $H$ of $\widehat{P}$ using \textsc{IncrementalHull}
  \State collect the facets of $H$ whose outward normal has negative last coordinate \Comment{lower hull}
  \State project those facets back to $\mathbb{R}^d$ by dropping the last coordinate
  \State \Return the projected complex
\EndFunction
\end{algorithmic}
\end{algorithm}

The orientation of the lower hull is decided by the sign of the last
coordinate of the outward normal, which is computed from the facet vertices
by the $d$-dimensional orientation predicate; this is the direct analogue of
the empty-circle test used in the plane.

\subsection{5. Parameter Selection}

For exact construction one computes, for each candidate $(d+1)$-simplex, the
sign of the determinant of the lifted coordinates, an empty-sphere test that
is the $d$-dimensional generalization of the incircle test. In floating
point the signs are unstable when the radius is much larger than the local
scale of the sites; symbolic perturbation of the coordinates resolves ties.
The relevant scale parameter is the minimum feature size of the site set,
which any robust implementation must estimate before choosing the working
precision.

\subsection{6. Complexity}

Lifting adds one dimension, so the hull bound
$O(n \log n + n^{\lfloor (d+1)/2 \rfloor})$ of
Theorem~\ref{thm:highdim_hull_randomized} becomes
$O(n \log n + n^{\lceil d/2 \rceil})$, matching
Theorem~\ref{thm:highdim_delaunay_complexity}. The Voronoi diagram is the
dual cell complex and is computed in the same time. The constants hidden in
the big-$O$ grow super-exponentially with $d$ (roughly $d^{d/2}$), so the
algorithm is practical only for $d \le 6$ or so; beyond that the output itself
is astronomically large and approximation replaces construction.

\subsection{7. Usages}

Delaunay triangulations in $\mathbb{R}^3$ are a standard tool in mesh
generation, and in dimension four they appear in the regular triangulations
used to compute volumes and mixed volumes. The lifting reduction is the
template for the more general concept of \index{regular
triangulation}regular triangulations, where the vertical shift of each lifted
point is allowed to vary, and it underlies the theory of secondary polytopes.
In statistical geometry, Delaunay complexes of point clouds appear as the
combinatorial part of persistent homology computations.

\subsection{8. Testing Methods and Edge Cases}

The primary correctness check is the empty-sphere predicate: every returned
simplex must admit a sphere through its vertices with no site strictly inside.
The dual check verifies that the projection of the lower hull has exactly the
combinatorial type of the Delaunay complex, which is provably the union of the
top faces of the so-called power diagram. Test configurations include
regularly spaced grids (maximal number of simplices), cocircular sites in a
hyperplane (flat lower hull), and the lifted cyclic polytope, which produces
the maximal Delaunay complex and stresses the manifold test. A hidden trap is
that the lower hull may project to a complex that is not a triangulation; the
implementation must not assume that every $d$-cell is a simplex.

\subsection{9. References}

The lifting construction and its properties are treated in \cite{berg2008}
for the plane and in \cite{matousek2002} for general dimension; the Delaunay
computations via hull are surveyed in \cite{deberg2008}. The randomized
incremental analysis underlying Theorem~\ref{thm:highdim_delaunay_complexity}
is in \cite{mulmuley1994}.

\section{Nearest-Neighbor Search}
\label{sec:highdim_nearest_neighbor}

\subsection{1. Overview}

Chapter~\ref{ch:proximity} solved planar nearest-neighbor queries with data
structures of linear size and polylogarithmic query time, and
Chapter~\ref{ch:range_searching} showed that range trees and $k$d-trees push
the same paradigm into higher dimensions, at the cost of
$O(n \log^{d-1} n)$ storage and $O(n^{1-1/d})$ query time. The degradation is
structural: for $d$ comparable to $\log n$ these bounds become polynomial and
the structure degenerates into a linear scan. The solution for growing
dimension is to give up exactness and to solve the
\index{approximate nearest neighbor}approximate nearest-neighbor problem, for
which sublinear query time is provably achievable by
\index{locality-sensitive hashing}locality-sensitive hashing (LSH). This
section develops the LSH framework of Indyk and Motwani, its instantiation on
Hamming and Euclidean spaces, and the near-optimal schemes built on top of it.

\subsection{2. Definitions}

\begin{definition}[Approximate near neighbor]
\label{def:highdim_ann}
Fix a metric space, a radius $r > 0$ and an approximation factor $c > 1$. The
$(r,c)$-\emph{near-neighbor} problem is: preprocess a set $P$ of $n$ points so
that, given a query $q$, if there exists a point $p \in P$ with
$\|p - q\| \le r$, the structure returns a point $p' \in P$ with
$\|p' - q\| \le c r$. The answer is allowed to be any point within distance
$c r$, and points within distance $r$ may be missed. The
\emph{approximate nearest-neighbor} problem asks for a point within factor
$c$ of the true nearest neighbor.
\end{definition}

\begin{definition}[Sensitive hash family]
\label{def:highdim_lsh}
A family $\mathcal{H}$ of hash functions on the point space is
$(r, c r, p_1, p_2)$-\emph{sensitive} for $p_1 > p_2$ if for uniformly random
$h \in \mathcal{H}$ and all points $p, q$,
\[
\Pr[h(p) = h(q)] \ge p_1 \quad\text{whenever}\quad \|p - q\| \le r,
\]
\[
\Pr[h(p) = h(q)] \le p_2 \quad\text{whenever}\quad \|p - q\| \ge c r.
\]
\end{definition}

The two probabilities are the workhorse parameters: close points collide
often, far points collide rarely.

\subsection{3. Theory}

The LSH construction amplifies a single weak hash family into a data
structure. Concatenate $k$ independent hash functions into
$g(p) = (h_1(p),\dots,h_k(p))$ and build $L$ independent hash tables under
independent concatenations $g_1,\dots,g_L$. A query hashes itself with each
$g_i$ and inspects the points in the resulting buckets. The two design
equations fix $k$ and $L$:

\[
k = \left\lceil \frac{\log n}{\log(1/p_2)} \right\rceil,
\qquad
L = n^{\rho},
\qquad
\rho = \frac{\log(1/p_1)}{\log(1/p_2)}.
\]

With these choices a fixed near neighbor $p^*$ collides with $q$ under a
single $g_i$ with probability $p_1^k \approx 1/L$, so it is found across the
$L$ tables with constant probability; meanwhile an arbitrary far point
collides with probability at most $p_2^k = 1/n$, so the expected number of far
candidates inspected in any one bucket is $1$.

\begin{theorem}[Indyk--Motwani bound]
\label{thm:highdim_lsh_bound}
Let $P \subset \{0,1\}^D$ be a set of $n$ points and let
$\mathcal{H}$ be the family of coordinate projections (each $h$ reads one
coordinate). Then the $(r, cr)$-near-neighbor problem on the Hamming cube can
be solved with a data structure of size $O(n^{1+\rho})$, query time
$O(D \cdot n^{\rho} \cdot \log n)$ and constant success probability, where
$\rho = \log(1/p_1)/\log(1/p_2)$ and $p_1 = 1 - r/D$, $p_2 = 1 - c r / D$.
For small $r/D$ the exponent is $\rho \approx 1/c$.
\end{theorem}

The Euclidean analogue uses the $p$-stable families of
Section~\ref{sec:highdim_random_projections}, which give
$\rho < 1/c$, and the asymptotically optimal schemes
reach $\rho = 1/c^2 + o(1)$. An exact nearest-neighbor query is reduced to
$(r,c)$-near-neighbor queries by a distance-doubling search: guess the
distance to the nearest neighbor as powers of two and run the near-neighbor
query for each guess.

\subsection{4. Pseudocode}

Algorithm~\ref{alg:highdim_lsh} is the preprocessing and query template. The
only free ingredients are the underlying family $\mathcal{H}$ and the
parameter pair $(k, L)$.

\begin{algorithm}[htbp]
\caption{Locality-Sensitive Hashing for $(r,c)$-Near Neighbor}
\label{alg:highdim_lsh}
\begin{algorithmic}[1]
\Require Points $P$, radius $r$, factor $c > 1$, and an
$(r, cr, p_1, p_2)$-sensitive family $\mathcal{H}$.
\Ensure A structure returning a point within distance $cr$ of $q$ if one is within $r$.
\Function{LSHPreprocess}{$P$}
  \State $k \gets \lceil \log_{1/p_2} n \rceil$ \Comment{concatenation length}
  \State $\rho \gets \log(1/p_1) / \log(1/p_2)$
  \State $L \gets n^{\rho}$ \Comment{number of tables}
  \For{$i = 1, \dots, L$}
    \State draw $h_{i,1}, \dots, h_{i,k}$ from $\mathcal{H}$
    \State let $g_i(p) = (h_{i,1}(p), \dots, h_{i,k}(p))$
    \State store every $p \in P$ in table $i$ under key $g_i(p)$
  \EndFor
\EndFunction
\Function{LSHQuery}{$q$}
  \State $seen \gets 0$
  \For{$i = 1, \dots, L$}
    \For{each point $p$ in bucket $g_i(q)$}
      \State $seen \gets seen + 1$
      \If{$\|p - q\| \le c r$} \State \Return $p$ \EndIf
      \If{$seen > 3 L$} \State \Return \textbf{none} \EndIf
    \EndFor
  \EndFor
  \State \Return \textbf{none}
\EndFunction
\end{algorithmic}
\end{algorithm}

\subsection{5. Parameter Selection}

The family $\mathcal{H}$ must be chosen for the metric at hand. On the Hamming
cube $\{0,1\}^D$ take coordinate projections; a point at distance $u$ collides
with the query with probability $1 - u/D$, yielding
$p_1 = 1 - r/D$ and $p_2 = 1 - c r/D$. For Euclidean space the $p$-stable
families of Section~\ref{sec:highdim_random_projections} hash by
$\lfloor (a \cdot x + b)/W \rfloor$ with a $p$-stable random vector $a$; the
bucket width $W$ trades the two probabilities. Increasing $k$ decreases both
probabilities but sharpens the ratio; increasing $L$ (equivalently increasing
$\rho$) raises the success probability and the storage. In practice one
chooses $W$ by a small grid search over a holdout set and raises $L$ until the
empirical success probability is acceptable. The hash tables themselves are
maintained as arrays indexed by $g_i$; a secondary collision hash on the
concatenated keys handles their variable size.

\subsection{6. Complexity}

Storage is $O(n \cdot (d + L)) = O(dn + n^{1+\rho})$. Each query evaluates
$L$ concatenations, each costing $O(dk)$ for the projections (or $O(k)$ per
non-zero coordinate on sparse data), and inspects at most $O(L)$ candidate
points, each verified in $O(d)$. The total query time is
$O(d \cdot L \cdot k + dL) = O(d n^{\rho} \log n)$. The trade-off
between the approximation factor and the exponent $\rho$ is the central
quantitative fact: $\rho \approx 1/c$ for the basic families, $\rho < 1/c$ for
the $p$-stable families, and $\rho = 1/c^2 + o(1)$ for the near-optimal
schemes. Lower bounds for hashing-based schemes state that $\rho \ge 1/c^2$ is
unavoidable, so the near-optimal schemes are optimal within the LSH framework.

\subsection{7. Usages}

LSH is the workhorse of high-dimensional similarity search in practice: near
duplicate detection in web corpora, image and video retrieval, de-duplication
in databases, and genome sequence search. It is also the standard subroutine
inside $k$-means-like clustering over high-dimensional data, where exact
nearest-neighbor structures are useless. Section~\ref{sec:highdim_ml} returns
to the machine learning applications; the connection to
\index{data structures}data structures for high-dimensional data is treated in
the handbook \cite{mehta2006}.

\subsection{8. Testing Methods and Edge Cases}

The canonical test compares the structure against a brute-force nearest
neighbor on small random data, measuring both the success probability and the
observed approximation factor over many queries. Edge cases: queries far from
all data (the structure must terminate, bounded by the $3L$ cutoff); clusters
of nearly coincident points (collision probability near $1$ makes buckets
large); and the parameter extremes $c \to 1^+$, where $\rho \to 1$, i.e. the
structure degenerates to a linear scan. A robust implementation must also
tolerate adversarial hash keys, where a secondary exact hash function
collides, by comparing candidate points with the true distance rather than by
trusting the bucket membership.

\subsection{9. References}

The framework and the Hamming analysis are due to \cite{indyk1998} and
\cite{gionis1999}. The Euclidean instantiation via $p$-stable distributions is
\cite{datar2004}, and the near-optimal schemes with $\rho = 1/c^2 + o(1)$ are
\cite{andoni2006}. The distance-concentration phenomenon that motivates
approximation is analyzed in \cite{beyer1999}. The modern practical HNSW
graphs, which combine the LSH idea with greedy navigation, are presented in
\cite{malkov2018}.

\section{Dimensionality Reduction}
\label{sec:highdim_dimensionality_reduction}

\subsection{1. Overview}

The recurring obstacle in Sections~\ref{sec:highdim_hulls} and
\ref{sec:highdim_voronoi_delaunay} is that combinatorial objects in
$\mathbb{R}^d$ have complexity exponential in $d$. The response of the field
is \index{dimensionality reduction}\emph{dimensionality reduction}: replace
the point set by a small set of sketches in a low-dimensional space, chosen so
that the quantities the application actually cares about (pairwise distances,
nearest neighbors, dot products) are approximately preserved. The guiding
statement is the \index{Johnson--Lindenstrauss lemma}Johnson--Lindenstrauss
lemma: $n$ points can be embedded into $\mathbb{R}^k$ with
$k = O(\log n)$ while preserving all pairwise distances within a factor
$1 \pm \varepsilon$. This section states the lemma, discusses the
information-theoretic intuition behind its logarithmic bound, and positions it
relative to the exact theory of metric embeddings and to the linear algebraic
methods used in practice.

\subsection{2. Definitions}

\begin{definition}[Distortion of an embedding]
\label{def:highdim_distortion}
A map $f: \mathbb{R}^d \to \mathbb{R}^k$ has \emph{distortion} at most
$1 + \varepsilon$ on a set $P$ if for all $p, q \in P$,
\[
(1 - \varepsilon)\,\|p - q\|^2 \;\le\; \|f(p) - f(q)\|^2 \;\le\; (1 + \varepsilon)\,\|p - q\|^2.
\]
The squared normalization is a convention that removes square roots from the
analysis; the equivalent statement for the un-squared distance holds with
$(1 \pm \varepsilon)$ replaced by $(1 \pm 3\varepsilon)$.
\end{definition}

A linear embedding is a matrix $A \in \mathbb{R}^{k \times d}$; it is
\emph{norm-preserving} for a unit vector $x$ if
$\|A x\|^2$ is close to $1$. Random constructions draw $A$ entrywise from a
zero-mean distribution scaled so each entry has variance $1/k$, in which case
$\mathbb{E}[\|A x\|^2] = \|x\|^2$ for every $x$: the random map is unbiased in
expectation, and the question becomes how tightly it concentrates.

\subsection{3. Theory}

Why is the target dimension logarithmic in $n$ and quadratic in
$1/\varepsilon$? The lower bound is forced by two simple arguments. First, $n$
equidistant points form a regular simplex in $\mathbb{R}^{n-1}$, and embedding
it faithfully needs roughly $\log n$ dimensions once the distances may only
be approximate. Second, the variance of the squared norm of a Gaussian
projection of a unit vector is about $2/k$, so guaranteeing a $1\pm\varepsilon$
multiplicative bound requires $k = \Omega(1/\varepsilon^2)$; concentration
cannot do better. Both bounds are achieved, up to constants, by the Gaussian
and Rademacher constructions of Section~\ref{sec:highdim_random_projections}.

\begin{theorem}[Johnson--Lindenstrauss]
\label{thm:highdim_jl_bound}
Let $P$ be a set of $n$ points in $\mathbb{R}^d$ and let
$\varepsilon \in (0,1)$. For any
$k \ge c\,\varepsilon^{-2} \log n$ with an absolute constant $c$, a random
$k \times d$ matrix with independent Gaussian or Rademacher entries of
variance $1/k$ satisfies the distortion bound of
Definition~\ref{def:highdim_distortion} on $P$ with probability at least
$1 - 1/n$.
\end{theorem}

The dimension reduction view extends from points to subspaces. A random
matrix with $k = O(d/\varepsilon^2)$ rows is a
\index{subspace embedding}subspace embedding: it preserves, up to factor
$1 \pm \varepsilon$, the norm of every vector lying in a fixed
$d$-dimensional subspace. This single statement packages the whole of the
modern randomized linear algebra (approximate SVD, least-squares regression,
low-rank approximation), where the subspace is the one spanned by the data
columns and the random map compresses the problem before a deterministic
solver runs on the sketch.

\subsection{4. Pseudocode}

The embedding itself is a single matrix multiplication
(Algorithm~\ref{alg:highdim_jl}); the data structure of this section is the
selection of the matrix and the dimension, discussed next.

\begin{algorithm}[htbp]
\caption{Johnson--Lindenstrauss Embedding}
\label{alg:highdim_jl}
\begin{algorithmic}[1]
\Require Points $p_1, \dots, p_n \in \mathbb{R}^d$, distortion $\varepsilon \in (0,1)$.
\Ensure Sketches $q_1, \dots, q_n \in \mathbb{R}^k$ with $k = O(\varepsilon^{-2} \log n)$.
\Function{JLEmbed}{$p_1, \dots, p_n$, $\varepsilon$}
  \State $k \gets \lceil c \ln n / \varepsilon^2 \rceil$ \Comment{the constant from the theory}
  \State draw a $k \times d$ matrix $A$ with i.i.d.\ entries of variance $1/k$
  \For{$i = 1, \dots, n$}
    \State $q_i \gets A p_i$
  \EndFor
  \State \Return $q_1, \dots, q_n$
\EndFunction
\end{algorithmic}
\end{algorithm}

The reduction in dimension is dramatic when the data is intrinsically
low-rank: computing pairwise distances in $\mathbb{R}^k$ with
$k \approx \log n$ instead of $\mathbb{R}^d$ with $d$ in the thousands is the
difference between feasible and infeasible. The cost is a single
$k \times d$ matrix multiplication per point, which on sparse data reduces to
a sum over non-zero entries.

\subsection{5. Parameter Selection}

The two parameters are the distortion $\varepsilon$ and the constant
$c$ hidden in the bound. For most applications $\varepsilon \in (0.05, 0.2)$
is the practical range: smaller values push $k$ above $\approx 100$ even for
moderate $n$, and larger values distort the metric enough to corrupt nearest
neighbor or clustering results. The constant $c$ is implementation dependent;
choosing $k = 2\,\varepsilon^{-2}\ln n$ is the common heuristic and works in
practice because the union bound over pairs is loose. When the downstream task
is clustering or $k$-NN classification, the distortion can often be larger
than the theory suggests, since only the local metric needs preservation.

\subsection{6. Complexity}

Embedding $n$ points costs $O(n d k)$ time with a dense matrix and
$O(\operatorname{nnz} \cdot k)$ on sparse data, where
$\operatorname{nnz}$ is the number of non-zeros of the data matrix. The
resulting sketch space is $O(n k) = O(n \varepsilon^{-2} \log n)$ words. A
random projection can never reduce the intrinsic rank of the data, but it does
compress the ambient dimension to the information-theoretic minimum
$k = \Theta(\log n)$ for distance preservation. When a deterministic
guarantee is required, one may verify the distortion empirically on the pair
set and rerun with fresh randomness until the bound holds.

\subsection{7. Usages}

Dimensionality reduction is the standard preprocessing in machine learning
pipelines: before running a $k$-means clustering, a $k$-NN classifier, or a
Gaussian kernel method, the data is embedded into $O(\log n)$ dimensions with
negligible loss (Section~\ref{sec:highdim_ml}). It is also the tool that makes
high-dimensional range searching and nearest-neighbor structures
(Sections~\ref{sec:highdim_nearest_neighbor} and
\ref{sec:range_search}) usable at scale, because the data structures of
Chapters~\ref{ch:range_searching} and \ref{ch:proximity} depend exponentially
on the ambient dimension.

\subsection{8. Testing Methods and Edge Cases}

Validate by computing the maximum pairwise distortion of the sketch against
the original over held-out pairs; the empirical maximum should be within
$1 \pm 2\varepsilon$. Edge cases: points already in low dimension (the map
should leave them essentially unchanged); nearly collinear clouds (the
sketched distances concentrate at the diameter, hiding the internal
structure); and points with a huge spread in norms, where multiplicative
error is acceptable but additive error is not. A hidden trap is numerical
noise: entries of variance $1/k$ are small when $k$ is large, and single
precision can destroy the sketch; use double precision in the accumulation.

\subsection{9. References}

The foundational statement is \cite{johnson1984}. The standard textbook
treatment, including the constant and the lower bounds, is
\cite{matousek2002}. Practical surveys and the connections to nearest
neighbor and learning appear in \cite{indyk1998} and \cite{harpeled2011}.

\section{Johnson--Lindenstrauss Lemma}
\label{sec:highdim_jl}

\subsection{1. Overview}

This section proves Theorem~\ref{thm:highdim_jl_bound}. The proof is a model
for the rest of high-dimensional analysis: it is a concentration argument
(one-dimensional, then union-bounded over pairs) wrapped around a linear map.
We give the two standard constructions --- Gaussian entries and Rademacher
entries --- and the concentration inequality behind each, and we discuss
tightness. The material of Section~\ref{sec:highdim_concentration} supplies
the probabilistic machinery, so this section may be read as its first
application.

\subsection{2. Definitions}

Let $A$ be a $k \times d$ matrix with i.i.d.\ entries $A_{ij}$ of mean $0$ and
variance $1/k$. For a fixed unit vector $x$, the vector $y = A x$ has
independent coordinates
\[
y_i = \sum_{j=1}^{d} A_{ij} x_j,
\]
each of mean $0$ and variance
$\sum_j (1/k) x_j^2 = 1/k$. Its squared norm
$\|y\|^2 = \sum_i y_i^2$ has expectation $1$; the whole proof is the statement
that $\|y\|^2$ is concentrated about $1$. For Gaussian entries, $k \|y\|^2$ is
a chi-squared random variable with $k$ degrees of freedom. For Rademacher
entries, a Hoeffding-type bound on $y_i$ followed by a sub-exponential tail on
$y_i^2$ gives the same result.

\subsection{3. Theory}

\begin{theorem}[One-dimensional concentration]
\label{thm:highdim_jl_vector}
Let $y = A x$ for a fixed unit vector $x$ and a $k \times d$ matrix with
independent Gaussian or Rademacher entries of variance $1/k$. For any
$t \in (0, 1)$,
\[
\Pr\big[\,\|y\|^2 \ge 1 + t\,\big] \le e^{-c k t^2},
\qquad
\Pr\big[\,\|y\|^2 \le 1 - t\,\big] \le e^{-c k t^2},
\]
for an absolute constant $c > 0$.
\end{theorem}

Proof of Theorem~\ref{thm:highdim_jl_vector}. For Gaussian entries, write
$y_i = g_i/\sqrt{k}$ with $g_i$ i.i.d.\ standard Gaussians, so
$k \|y\|^2 = \sum_i g_i^2 \sim \chi^2_k$. The chi-squared tail bound
\[
\Pr\left[ \frac{1}{k}\sum_{i=1}^{k} g_i^2 - 1 \ge t \right]
\le \exp\left( - \frac{k}{2}(t - \ln(1+t)) \right),
\]
together with $t - \ln(1+t) \ge t^2/4$ for $t \in (0,1)$, gives the upper
tail; the lower tail is symmetric by the same argument on $-y$. For Rademacher
entries, $y_i$ is a sum of independent bounded variables, so
 Hoeffding's inequality gives a Gaussian tail on $y_i$, and summing the
 truncated squares $y_i^2 \wedge M$ with a Bernstein bound yields the same
 exponential tail.


\begin{theorem}[Johnson--Lindenstrauss, proved]
\label{thm:highdim_jl_proved}
For $n \ge 2$, $\varepsilon \in (0,1)$, and
$k \ge 8 \varepsilon^{-2} \ln n$, the map of
Theorem~\ref{thm:highdim_jl_bound} satisfies the distortion bound on all
$\binom{n}{2}$ pairs simultaneously with probability at least $1 - 1/n$.
\end{theorem}

Proof. For a fixed pair $(p, q)$, apply Theorem~\ref{thm:highdim_jl_vector} to
the unit vector $x = (p - q)/\|p - q\|$, noting that
$A(p - q) = A p - A q$ and that the squared norm scales by
$\|p - q\|^2$. The two tails each cost at most $e^{-c k \varepsilon^2}$;
union bound over the $\binom{n}{2}$ pairs gives failure probability at most
$n^2 e^{-c k \varepsilon^2} \le 1/n$ for
$k \ge (3/c) \varepsilon^{-2} \ln n$. Absorbing the constants yields the
claimed bound.

\subsection{4. Pseudocode}

The construction is Algorithm~\ref{alg:highdim_jl}; its correctness follows
from the two theorems. The Rademacher variant replaces each Gaussian by
$\pm 1/\sqrt{k}$ with equal probability, which avoids random-number generation
of continuous distributions and enables the hashed implementation of
Section~\ref{sec:highdim_random_projections}.

\subsection{5. Parameter Selection}

The constant $8$ in the bound is an artifact of the union bound and the
concentration constant; smaller values work in practice. Two practical
refinements matter. First, centering: subtracting the mean point before
embedding does not change distances but reduces the scale against which
floating point error is measured. Second, when only $k$ can be chosen
(not the constant), setting $k = \lceil 2 \ln n / \varepsilon^2 \rceil$ is the
de facto standard; Section~\ref{sec:highdim_random_projections} discusses the
consequences of the choice of the distribution.

\subsection{6. Complexity}

The embedding is a single dense multiplication, $O(n d k)$ time and $O(dk)$
space for the matrix; with $k = O(\log n)$ this is $O(n d \log n / \varepsilon^2)$
for the whole data set. The Rademacher version has the same complexity but
faster constants, and the sparse variants of
Section~\ref{sec:highdim_random_projections} reduce the multiplication cost to
the number of non-zeros.

\subsection{7. Usages}

The lemma is invoked wherever a high-dimensional metric is to be preserved at
logarithmic cost: before any nearest-neighbor or clustering computation
(Section~\ref{sec:highdim_ml}), inside kernel methods via random features, and
as the standard tool in data stream sketching, where only the projection of
each incoming point is stored. It also underlies the theoretical reduction of
many high-dimensional problems to their $d = O(\log n)$ instances.

\subsection{8. Testing Methods and Edge Cases}

Verify the embedding empirically on random Gaussian point sets with
$n = 10^5$, $d = 10^3$, $\varepsilon = 0.1$: the maximum distortion should
stay below $1 + 0.1$ with high probability. Edge cases: nearly identical
points (their difference is a near-unit vector, still concentrated); points at
the origin (trivially preserved); adversarial points chosen after the matrix
is fixed (the union bound fails, and one must either resample or verify
deterministically). Numerical care is required: the entries of $A$ scale as
$1/\sqrt{k}$ and the products $A p$ accumulate round-off proportional to
$\|p\| \sqrt{k}$.

\subsection{9. References}

The lemma is from \cite{johnson1984}; the modern proof via concentration is
standard and is presented in \cite{matousek2002}. The chi-squared tail bound
used in the proof is derived in Section~\ref{sec:highdim_concentration}, and
the Rademacher and sparse constructions are \cite{achlioptas2003}. The
dimension bound is essentially optimal, as discussed in
\cite{matousek2002}.

\section{Random Projections}
\label{sec:highdim_random_projections}

\subsection{1. Overview}

A random projection is a concrete matrix distribution with the property that
every norm is approximately preserved. Section~\ref{sec:highdim_jl} proved
that Gaussian and Rademacher matrices work; this section is about the
engineering: which distributions are used in practice, how to hash them so
that a dense matrix is never materialized, how to use them for
\index{subspace embedding}subspace embeddings and approximate linear algebra,
and how the same idea instantiates the $p$-stable hash families of
Section~\ref{sec:highdim_nearest_neighbor}.

\subsection{2. Definitions}

\begin{definition}[Random projection matrix]
\label{def:highdim_random_projection}
A \emph{random projection} is a random $k \times d$ matrix $A$ whose entries
are independent, of mean $0$ and variance $1/k$, drawn from one of the
distributions below. Applying $A$ to a point set is
\emph{random projection}; the image points are called \emph{sketches}.
\end{definition}

The three classical entry distributions are:
\begin{itemize}
\item \emph{Gaussian}: $A_{ij} \sim \mathcal{N}(0, 1/k)$. Optimal constants,
costs one normal draw per entry.
\item \emph{Rademacher}: $A_{ij} \in \{+1/\sqrt{k}, -1/\sqrt{k}\}$ each with
probability $1/2$. No continuous sampling; supports hashing.
\item \emph{Sparse (Achlioptas)}: $A_{ij} \in \{0, \pm \sqrt{3/k}\}$ with
$\Pr[\pm] = 1/6$ and $\Pr[0] = 2/3$. Each entry has variance
$(3/k) \cdot (1/3) = 1/k$; only one third of the entries are non-zero.
\end{itemize}

A \emph{$p$-stable} random vector $a \in \mathbb{R}^d$ is one such that
$a \cdot x$ has the same distribution as $\|x\|_p \, X$ for a fixed
$p$-stable random variable $X$; Gaussians are $2$-stable, Cauchy variables are
$1$-stable.

\subsection{3. Theory}

The subspace embedding theorem packages the norm-preservation guarantee
uniformly over a subspace rather than over a finite set of points.

\begin{theorem}[Subspace embedding]
\label{thm:highdim_subspace_embedding}
Let $V \subseteq \mathbb{R}^d$ be a linear subspace of dimension $D$ and let
$A$ be a $k \times d$ matrix with independent Gaussian entries of variance
$1/k$ and $k = O(D / \varepsilon^2)$. Then with probability at least
$1 - \exp(-\Omega(D))$,
\[
(1 - \varepsilon)\, \|x\|^2 \;\le\; \|A x\|^2 \;\le\; (1 + \varepsilon)\, \|x\|^2
\qquad \text{for all } x \in V.
\]
\end{theorem}

The proof is an $\varepsilon$-net argument: a bounded net of
$O(2^{O(D)})$ points in the unit sphere of $V$ carries the norm of every
vector up to an additive error of $\varepsilon$, and the union bound over the
net absorbs the exponential dependence on $D$. The Gaussian distribution is
rotation-invariant, which is what makes it the right choice for the uniform
statement; the Rademacher version holds with slightly more rows.

The $p$-stable families are the bridge to LSH. Hash with
\[
h_{a,b}(x) = \left\lfloor \frac{a \cdot x + b}{W} \right\rfloor,
\qquad a \sim p\text{-stable},\; b \sim \mathrm{Uniform}[0, W].
\]
For two points at Euclidean distance $u$, the difference $a \cdot (x - y)$ is
$u X$ for a $p$-stable $X$, so the collision probability is
\[
p(u) = \mathbb{E}\left[ \left(1 - \frac{|u X|}{W}\right)_+ \right],
\]
a decreasing function of $u$ with $p(0) = 1$ and $p(u) \to 0$. Choosing
$W \approx r$ makes $p_1 = p(r)$ and $p_2 = p(c r)$ small together, which
yields $\rho < 1/c$ for the $2$-stable (Gaussian) case.

\subsection{4. Pseudocode}

Algorithm~\ref{alg:highdim_random_projection} gives the generic application
step; the distribution is a parameter.

\begin{algorithm}[htbp]
\caption{Random Projection}
\label{alg:highdim_random_projection}
\begin{algorithmic}[1]
\Require Points $p_1, \dots, p_n \in \mathbb{R}^d$, target dimension $k$.
\Ensure Sketches $q_1, \dots, q_n \in \mathbb{R}^k$.
\Function{RandomProject}{$P$, $k$}
  \State choose an entry distribution (Gaussian, Rademacher, or sparse Achlioptas)
  \State draw a $k \times d$ matrix $A$ from that distribution, scaled to variance $1/k$
  \For{$i = 1, \dots, n$}
    \State $q_i \gets A p_i$
  \EndFor
  \State \Return $q_1, \dots, q_n$
\EndFunction
\end{algorithmic}
\end{algorithm}

The hashed implementation avoids materializing $A$: instead of storing the
matrix, store for each row a pseudo-random hash of the input coordinate into
$\{+1, -1, 0\}$ (for the sparse family) together with the $p$-stable sample
value (for the Gaussian family), and accumulate $q_i$ coordinate by coordinate.
This reduces the storage from $O(dk)$ to $O(k)$ words.

\subsection{5. Parameter Selection}

The target dimension $k$ is set by the application: $\varepsilon^{-2}\log n$
for distance preservation, $O(D/\varepsilon^2)$ for a known subspace of
dimension $D$, or a fixed budget in streaming settings. The bucket width $W$
of the $p$-stable hash is chosen by grid search so that the ratio
$p_1/p_2$ is maximized at the operating radius $r$; in practice $W$ a few
times $r$ balances the collision probability against the bucket noise. The
sparse Achlioptas family is preferred over the dense Gaussian when
$d$ is large and the data is sparse, because the expected number of non-zero
products per output coordinate is $d/3$ rather than $d$.

\subsection{6. Complexity}

Dense application costs $O(n d k)$; the sparse family costs
$O(n \cdot \operatorname{nnz})$ regardless of $k$ when the hashed form is
used, since each non-zero input coordinate touches at most one entry per
output coordinate. The $p$-stable hash evaluates in $O(d)$ per query (one dot
product per table), which is why the query time of
Algorithm~\ref{alg:highdim_lsh} carries the factor $d$. The subspace embedding
compresses any $n \times d$ data matrix to $k \times d$ (or $n \times k$,
whichever is smaller) before the expensive linear algebra runs, changing the
dominant term of an SVD or least-squares solve from cubic in the ambient
dimension to cubic in the sketch dimension.

\subsection{7. Usages}

Random projections are used for approximate SVD and principal component
analysis, for least-squares regression on wide matrices, for $k$-means
clustering at scale, for kernel methods via random features, and for
similarity search on text and image data through the $p$-stable LSH families.
The sparse matrix products that appear when sketching Gram matrices are
sped up by the algorithms of \cite{yuster2004}.

\subsection{8. Testing Methods and Edge Cases}

Test that $\mathbb{E}[\|A x\|^2] = \|x\|^2$ on random unit vectors for each
distribution (a statistical check of the scaling) and that the empirical
distortion over the data set stays below $1 + \varepsilon$. Edge cases: the
zero vector (preserved exactly, do not divide by it in normalized checks);
points with non-zero mean (center first, as in
Section~\ref{sec:highdim_jl}); and the pathological case where the data
itself has an adversarial sparse structure aligned with the sampled
coordinates, which the hashed implementation can under-save. Whenever the
guarantee must be deterministic, run the map twice with independent randomness
and keep the better sketch.

\subsection{9. References}

The entry distributions are analyzed in \cite{achlioptas2003}; the
$p$-stable families are \cite{datar2004}; the subspace embedding view and its
applications to randomized linear algebra are developed in
\cite{harpeled2011}. The classical statements appear in \cite{johnson1984}
and \cite{matousek2002}.

\section{Concentration of Measure}
\label{sec:highdim_concentration}

The probabilistic machinery behind every random construction of this chapter
is the \index{concentration of measure}concentration of measure: in high
dimension, smooth functions of many independent (or independent-looking)
variables are essentially constant. This section develops the phenomenon in
its geometric forms --- the Gaussian annulus, the volume of balls, Levy's
lemma on the sphere --- and then shows how the same principle rescues
algorithms whose worst-case analysis is catastrophic: the average-case and
smoothed behavior of the simplex method, and the sampling-based computation of
volumes of convex bodies.

\begin{theorem}[Gaussian annulus]
\label{thm:highdim_gaussian_annulus}
Let $g \sim \mathcal{N}(0, I_d)$. Then with probability at least
$1 - 2 e^{-c d \varepsilon^2}$,
\[
(1 - \varepsilon)\sqrt{d} \;\le\; \|g\| \;\le\; (1 + \varepsilon)\sqrt{d},
\]
for an absolute constant $c > 0$.
\end{theorem}

The theorem says that a high-dimensional Gaussian vector is almost surely at
distance $\sqrt{d}$ from the origin, within a multiplicative error that
shrinks with $d$. The intuition is a variance computation: each coordinate
contributes variance $1$ to $\|g\|^2$, which has expectation $d$ and standard
deviation $\sqrt{2d}$, so the relative deviation $\sqrt{2d}/d = \sqrt{2/d}$
tends to zero. The proof is the chi-squared tail of
Section~\ref{sec:highdim_jl}: $\|g\|^2$ is a chi-squared variable with $d$
degrees of freedom, and
\[
\Pr\left[ \|g\|^2 \ge (1+t)d \right] \le e^{-c d t^2}
\]
follows from the exponential Markov inequality applied to
$e^{\lambda g^2}$. This is exactly the one-dimensional statement
(Theorem~\ref{thm:highdim_jl_vector}) that powered the
Johnson--Lindenstrauss lemma, in its pure form.

The same phenomenon concentrates the volume of the unit ball. Writing
\[
\operatorname{vol}(B^d) = \frac{\pi^{d/2}}{\Gamma(d/2 + 1)},
\]
the fraction of the volume lying within distance $\delta$ of the boundary of
the unit ball is
\[
1 - (1 - \delta)^d,
\]
which tends to $1$ exponentially fast. Equivalently, the ball of radius
$1 - \delta$ contains an exponentially small fraction of the volume: high
dimensions are empty in their interior and populated on their boundary shell.
This is why a uniform grid with spacing $\varepsilon$ in the unit cube needs
$\varepsilon^{-d}$ points --- the volume-sampling version of the curse of
dimensionality (Section~\ref{sec:highdim_curse}).

\begin{theorem}[Levy's lemma]
\label{thm:highdim_levy}
Let $f: S^{d-1} \to \mathbb{R}$ be $L$-Lipschitz with respect to geodesic
distance on the unit sphere. Then
\[
\Pr\left[ |f - \mathbb{E}[f]| \ge t \right] \le 2 \exp\left( - \frac{c d t^2}{L^2} \right)
\]
for an absolute constant $c > 0$.
\end{theorem}

Levy's lemma is the general concentration statement for the sphere: every
Lipschitz function is nearly constant on most of the sphere. The proof passes
through the isoperimetric inequality on the sphere, which says that among all
sets of a given measure the caps minimize the boundary; integrating the
isoperimetric profile along the level sets of $f$ yields the exponential
tail. The two corollaries most used in this chapter are the Gaussian annulus
and the concentration of Lipschitz functions on the sphere, which together
imply the one-dimensional concentration of
Theorem~\ref{thm:highdim_jl_vector} and hence the Johnson--Lindenstrauss
lemma. A further corollary is that a random point on the sphere is within
$O(1/\sqrt{d})$ of any fixed equator with overwhelming probability ---
high-dimensional spheres are concentrated near every great circle, a fact
with no analogue in the plane.

\begin{remark}
Concentration is a double-edged sword. It makes random maps and random
estimators reliable, which is what makes the algorithms of this chapter
work, but it is also exactly the force that destroys nearest neighbors
(Section~\ref{sec:highdim_curse}): when all points of a Gaussian cloud are at
distance $\sqrt{d}$ from one another up to a factor $1 + O(1/\sqrt{d})$, the
distinction between "nearest" and "far" collapses.
\end{remark}

\subsection{Sampling convex bodies}

The same concentration phenomenon, in its algorithmic form, is the engine of
modern convex geometry. Consider the problem of approximating the volume of a
convex body given only by a membership oracle --- the generalization of the
polyhedron volume computations of Chapter~\ref{ch:convex_polytopes} to
implicit bodies.

\begin{theorem}[Approximate volume of convex bodies]
\label{thm:highdim_volume}
There is a randomized algorithm that, given a membership oracle for a convex
body $K \subset \mathbb{R}^d$ of positive volume and an approximation factor
$1 + \varepsilon$, outputs in time polynomial in $d$, $1/\varepsilon$, and
the bit complexity of the input a value within a factor $1 \pm \varepsilon$ of
$\operatorname{vol}(K)$, with high probability. Computing the volume exactly
is \#P-hard.
\end{theorem}

The algorithm is the paradigm of randomized geometric integration. It samples
nearly uniform points from $K$ using a random walk --- the hit-and-run
walk, which moves from the current point along a random direction to a random
point inside $K$ on that chord --- and it combines the samples with a
multiphase decomposition: $K$ is intersected with a growing sequence of balls
of radius $r_1 < r_2 < \dots < r_m = R$ such that each shell has bounded
volume ratio, and the volume ratios are estimated by counting samples in the
successive shells. The product of the ratios telescopes to
$\operatorname{vol}(K)$, and the concentration of the counts around their
means is the only error source. The key analytic fact is that the hit-and-run
walk mixes in $O(d^3)$ steps from a warm start --- polynomial in the
dimension, in stark contrast with the naive grid, which needs
$\varepsilon^{-d}$ points. This is concentration working in favor of the
algorithm: it makes the sampled estimators reliable.

The same sampling machinery computes integrals over convex bodies and,
via rejection, the probabilities of events in convex sets; it is the geometric
core of the Monte Carlo methods in Bayesian statistics and of the counting
problems reducible to volume computation.

\subsection{Average case and smoothed analysis of the simplex method}

The linear programming algorithms of Chapter~\ref{ch:linear_programming} are
polynomial for fixed dimension but the simplex method has exponential
worst-case behavior on the Klee--Minty cubes. High-dimensional geometry
explains why the simplex method works anyway in practice: the average case is
benign because the combinatorial objects that would make it slow --- long
chains of pivots --- are exponentially unlikely under natural random models.

\begin{theorem}[Borgwardt's average-case bound]
\label{thm:highdim_borgwardt}
Under a rotationally symmetric distribution on the coefficients, the simplex
method performs an expected number of pivot steps that is polynomial in $n$
and $d$.
\end{theorem}

Borgwardt's analysis shows that a random linear program has, with high
probability, a shadow-vertex pivot path of polynomial length: projecting the
feasible polyhedron onto the plane spanned by the objective and a fixed
direction yields a polygon whose size is the number of pivots, and the
spherical concentration of the coefficients keeps that projected polygon
small. The geometry is pure concentration of measure on the sphere: the
relevant pivots are those whose normal vectors lie in a thin band, and the
probability that the band contains many of them is exponentially small.

\begin{theorem}[Smoothed analysis of the simplex method]
\label{thm:highdim_simplex_smoothed}
Let a linear program be obtained by perturbing an arbitrary instance with
independent Gaussian noise of standard deviation $\sigma$. The expected number
of simplex pivots (under the shadow-vertex rule) is polynomial in $n$, $d$,
and $1/\sigma$.
\end{theorem}

Spielman--Teng's smoothed analysis is the definitive resolution of the
paradox: no instance is hard in the worst case once it is perturbed, and the
perturbation need be only polynomially small. The proof analyzes the expected
number of changes of the optimal basis under perturbation and again exploits
that the expected number of vertices of a random projection of a polytope is
polynomial. The simplex method and its smoothed behavior are treated in
detail in Chapter~\ref{ch:linear_programming}; the point here is that
concentration, which wrecks nearest-neighbor search, is the very mechanism
that makes the simplex method fast. The worst-case exponential examples are
artifacts of a negligible set of adversarial inputs.

\begin{remark}
The three theorems of this section --- Gaussian concentration, polynomial
volume computation, and polynomial smoothed pivoting --- share a single
methodological message: in high dimension, the space of inputs is dominated by
typical inputs, and typical geometry is tractable. Algorithm design for high
dimensions must therefore be \emph{average-case} or \emph{smoothed}, or it
must be \emph{approximate} with a randomized guarantee, as in
Sections~\ref{sec:highdim_nearest_neighbor} and \ref{sec:highdim_jl}.
\end{remark}

The references for this section are \cite{matousek2002} (Gaussian annulus,
Levy's lemma, volumes), \cite{dyerfriezekannan1991} (volume computation),
\cite{borgwardt1987} (average case of the simplex method), and
\cite{spielteng2004} (smoothed analysis).

\section{Curse of Dimensionality}
\label{sec:highdim_curse}

\begin{definition}[Curse of dimensionality]
\label{def:highdim_curse}
The \index{curse of dimensionality}\emph{curse of dimensionality} is the
collection of phenomena, observed in Section~\ref{sec:highdim_concentration},
in which the running time or storage of an algorithm, or the sample complexity
of a statistical estimator, grows like a quantity exponential in the
dimension --- $n^{\lfloor d/2 \rfloor}$, $\varepsilon^{-d}$, $2^d$, or
$\exp(c d)$ --- even when the dependence on the number of points $n$ is
benign. It manifests in three distinct forms for this book: combinatorial
(complexity of geometric structures), volumetric (the grid bound), and
metric (the collapse of distance information).
\end{definition}

\subsection{The combinatorial curse}

The combinatorial form is Theorem~\ref{thm:highdim_hull_complexity}: hulls,
Voronoi diagrams, and Delaunay complexes have
$n^{\Theta(d/2)}$ cells. This is not an artifact of the specific structure;
it is the number of cells of any arrangement of $n$ hyperplanes in
$\mathbb{R}^d$, which is $\Theta(n^d)$, and the number of facets of the
polar polytope, which is the same count. The complexity is intrinsic: even
just listing the Delaunay complex of a cyclic-polytope sample requires
$n^{\lceil d/2 \rceil}$ output, so no exact data structure in high dimension
can beat a linear scan on the input of its own output. Section
\ref{sec:highdim_voronoi_delaunay} and Section~\ref{sec:highdim_hulls} pay the
price in their complexity analyses; here we note that the bound is optimal
and that approximation is the only escape.

The same combinatorial explosion governs the data structures of
Chapter~\ref{ch:range_searching}: a $k$d-tree answering orthogonal range
queries in time $O(n^{1-1/d} + k)$ degrades to $O(n)$ as soon as $d$ grows,
and a range tree storing $n \log^{d-1} n$ points is pointless when $d \ge
\log n$. Section~\ref{sec:rangesearch_curse} documented the transition, and
Section~\ref{sec:rangesearch_simplex} proved that the transition is
unavoidable: any data structure with near-linear space must spend
$\Omega(n^{1-1/d})$ time per simplex range query in the worst case. In high
dimension, $n^{1-1/d}$ is $n$; range searching in high dimensions is a linear
scan.

\subsection{The volumetric curse}

The volumetric form is the grid bound derived in
Section~\ref{sec:highdim_concentration}: covering the unit cube with boxes of
side $\varepsilon$ needs $\varepsilon^{-d}$ boxes, and the fraction of the
volume of the unit ball within distance $\delta$ of the boundary is
$1 - (1-\delta)^d$. Every algorithm that partitions space must pay this cost,
because the volume itself is concentrated near the surface. This is the 
mechanism behind the $\varepsilon^{-d}$ sample complexity of nonparametric
statistics and behind the failure of uniform grid methods, and it is why
Section~\ref{sec:highdim_dimensionality_reduction} embeds before
partitioning. The dimensionality reduction escape works because the
\emph{intrinsic} dimension of the data (the dimension of the manifold or
subspace it occupies) is typically far smaller than the ambient dimension;
the grid bound is a bound on the ambient dimension alone.

\subsection{The metric curse}

The metric form is the most damaging for machine learning, because it attacks
the very meaning of similarity.

\begin{theorem}[Distance concentration]
\label{thm:highdim_distance_concentration}
Let $P$ be a set of $n$ i.i.d.\ points from a distribution supported in
$\mathbb{R}^d$, and let $q$ be an independent point. Under mild moment and
concentration conditions (satisfied by Gaussian data), for every fixed
$n$ the ratio
\[
\operatorname{rel}(q) = \frac{\min_{p \in P} \|p - q\|}{\max_{p \in P} \|p - q\|}
\]
converges to $1$ in probability as $d \to \infty$.
\end{theorem}

The statement is the analysis of the phenomenon noticed in
Section~\ref{sec:highdim_concentration}: all pairwise distances are within a
relative window of width $O(1/\sqrt{d})$ around a common scale, so the
nearest neighbor is only marginally closer than the farthest point. A
nearest-neighbor classifier whose discrimination rests on the gap between
nearest and second-nearest loses its meaning: there is no signal left.
The practical consequence is not that high-dimensional nearest-neighbor search
is impossible --- Section~\ref{sec:highdim_nearest_neighbor} shows sublinear
queries are achievable --- but that the \emph{queries} must be approximate and
that the \emph{data} must first be embedded so that the effective dimension of
the distance is small. The results are due to \cite{beyer1999}, who also give
the precise conditions under which the relative contrast degenerates, and the
phenomenon is discussed further in Section~\ref{sec:highdim_ml}.

\subsection{Hardness of worst-case high-dimensional problems}

Beyond the geometric and metric barriers sit the computational ones. The
worst-case version of exact nearest-neighbor search is as hard as brute force:
in a cell-probe model, any data structure answering exact queries in the
Hamming cube requires near-linear query time in the worst case, so the
approximation guarantee of Section~\ref{sec:highdim_nearest_neighbor} is not
a convenience but a necessity. More broadly, the PCP framework --- the
\index{probabilistically checkable proofs}PCP machinery developed to prove
hardness of approximation --- established that many optimization problems over
high-dimensional inputs are hard to approximate even approximately, in the
worst case. The connection runs through the same high-dimensional geometry:
a proof-checker verifies a witness by reading a few randomly chosen
coordinates, and the soundness of the scheme rests on the concentration of
the acceptance probability over the random coordinate choices. The
foundational references are \cite{arora1994} and the survey in
\cite{harpeled2011}. The message for algorithm design is the same as in
Section~\ref{sec:highdim_concentration}: the tractable regime in high
dimensions is the typical case and the approximate case, not the worst exact
case.

\subsection{The escapes}

Three complementary escapes from the curse run through this chapter:
\begin{itemize}
\item \emph{Dimensionality reduction} (Sections
\ref{sec:highdim_dimensionality_reduction}, \ref{sec:highdim_jl},
\ref{sec:highdim_random_projections}): reduce the ambient dimension to the
intrinsic one; the cost is a small relative distortion, which many downstream
tasks tolerate.
\item \emph{Approximate and randomized queries} (Section
\ref{sec:highdim_nearest_neighbor}): replace the exact answer by one that is
correct up to a factor $c$ with constant probability; the query cost drops from
$n$ to $n^{\rho}$ with $\rho < 1$.
\item \emph{Data-dependent and worst-case-intrinsic algorithms}: exploit the
\index{intrinsic dimension}intrinsic dimension or the 
\index{smoothed analysis}smoothed complexity of the instance
(Section~\ref{sec:highdim_concentration}), where the typical behavior is
polynomial even though the worst case is exponential.
\end{itemize}

\section{Applications to Machine Learning}
\label{sec:highdim_ml}

The pipeline of modern machine learning --- a high-dimensional feature vector,
a similarity or kernel, and an optimization or search step --- is, from the
point of view of this chapter, a high-dimensional computational geometry
problem with $n$ points in $\mathbb{R}^d$, and nearly every technique of the
preceding sections appears in it. This section walks through the pipeline and
maps each stage to the geometric tool. The statistical and algorithmic
foundations of the geometric view of learning are developed in
Chapter~\ref{ch:geometry_data_ml} and in the approximation chapter's
treatment of high-dimensional algorithms
(Section~\ref{sec:approx_high_dimensional}).

\subsection{$k$-nearest-neighbor classifiers}

The $k$-nearest-neighbor classifier labels a query by the majority label among
its $k$ nearest neighbors in the training set. Its accuracy depends
completely on the geometry of the metric: if the nearest neighbor is not
meaningfully closer than the rest, the classifier degenerates to the majority
class. The distance-concentration analysis of Section~\ref{sec:highdim_curse}
(\cite{beyer1999}) makes the failure precise: on high-dimensional data with no
intrinsic structure, the relative contrast between nearest and farthest points
tends to $1$, and the classifier has no signal. The remedies are the escapes
of Section~\ref{sec:highdim_curse}: embed the data first
(Section~\ref{sec:highdim_dimensionality_reduction}), or use an approximate
nearest-neighbor structure (Section~\ref{sec:highdim_nearest_neighbor}).
This is why, in practice, high-dimensional $k$-NN pipelines run a random
projection (or a learned linear map) before the search: the projection is not
a loss of information but a denoising of the metric. The state of the art in
production systems combines these ideas in the graph-based HNSW structures of
\cite{malkov2018}, whose greedy navigation over a proximity graph replaces the
exact structures that the curse kills.

\subsection{Dimensionality reduction as preprocessing}

The same argument applies to clustering, outlier detection, and kernel
methods. Before a Gaussian-kernel SVM on $d$-dimensional features, one can
replace the exact kernel by a random-feature map: draw a random projection
$A$ and use the finite-dimensional feature map
$\phi(x) = (\cos, \sin)$-features derived from $A x$, which by the 
Johnson--Lindenstrauss lemma approximates the kernel's geometry
(Section~\ref{sec:highdim_random_projections}). This is a direct application
of the lemmas of Sections~\ref{sec:highdim_jl} and
\ref{sec:highdim_random_projections} and is one of the few places where the
concentration theorems pay for themselves twice: once for the map's
faithfulness and once for the kernel's approximation. The data management side
of these pipelines --- indexing the sketches, hashing the features, and
maintaining the tables --- is surveyed in the handbook \cite{mehta2006}.

\subsection{Sparse data and fast linear algebra}

Real feature vectors are sparse: each example has a handful of non-zero
coordinates out of millions. Sparse random projections
(Section~\ref{sec:highdim_random_projections}) exploit this directly, and the
matrix multiplications that build Gram matrices and covariance estimates from
sparse inputs are accelerated by the sparse matrix product algorithms of
\cite{yuster2004}. The interplay between sparsity and geometry is itself a
high-dimensional phenomenon: the concentration results of
Section~\ref{sec:highdim_concentration} govern how many samples are needed to
estimate inner products and distances from sketches, and the same machinery
bounds the statistical error of the random-feature approximations.

\subsection{The geometric picture of learning}

The geometric reading of learning theory is the one this chapter makes
natural. A data set is a point cloud; the tasks are geometric queries
(nearest neighbors, kernels, clustering, projections); the sample complexity
of a hypothesis class is governed by the combinatorial complexity of the
geometry --- the VC dimension of ranges in $\mathbb{R}^d$ is $O(d)$ and
$\varepsilon$-nets of size $O(\frac{d}{\varepsilon}\log\frac{1}{\varepsilon})$
exist (Section~\ref{sec:rangesearch_kd_trees} and the related material in
Chapter~\ref{ch:range_searching}), which is the polynomial-in-$d$ side of the
curse. The exponential side is the sample complexity of volume-based
estimators (Section~\ref{sec:highdim_concentration}). Learning in high
dimensions is possible exactly when the data has low intrinsic dimension, and
dimensionality reduction is the algorithmic expression of that belief. The
standard reference for the geometric viewpoint is
Chapter~\ref{ch:geometry_data_ml}; the approximation algorithms that make
learning tractable in high dimensions are covered in
Section~\ref{sec:approx_high_dimensional} of the approximation chapter.

\begin{remark}
The most common mistake in high-dimensional practice is to assume that
Euclidean distance on raw features is a meaningful measure of similarity. The
analysis of Section~\ref{sec:highdim_curse} shows when it is not, and the
techniques of this chapter show how to fix it: normalize, embed, and
approximate. Each step is justified by a theorem in this chapter, and each
step is cheap enough for production systems.
\end{remark}

\section{Exercises}
\label{sec:highdim_exercises}

\begin{enumerate}
\item Prove that computing the convex hull of $n$ points in $\mathbb{R}^d$
requires $\Omega(n \log n)$ time even when the hull has $O(n)$ facets, by
reducing sorting to a hull computation in a plane contained in
$\mathbb{R}^d$.
\item Show that the expected number of facets created by the
randomized incremental hull algorithm (Algorithm~\ref{alg:highdim_hull}) on a
random sample is $O(n^{\lfloor d/2 \rfloor})$, and derive the running time of
Theorem~\ref{thm:highdim_hull_randomized} from the Clarkson--Shor charging
argument.
\item Prove the lifting statement of Section~\ref{sec:highdim_voronoi_delaunay}
in full generality: a sphere in $\mathbb{R}^d$ with center $c$ and radius $r$
corresponds to a hyperplane through the lifted points, and a point lies inside
the sphere iff its lift lies below the hyperplane. Conclude the empty-sphere
test from the orientation of the lifted simplex.
\item Let $\mathcal{H}$ be the family of coordinate projections on
$\{0,1\}^D$. Compute $p_1$, $p_2$, and $\rho$ for the $(r, cr)$-near-neighbor
problem and verify the exponent $\rho \approx 1/c$ for small $r/D$ in
Theorem~\ref{thm:highdim_lsh_bound}. Show that the query of
Algorithm~\ref{alg:highdim_lsh} inspects $O(L)$ points in expectation and
terminates with constant success probability.
\item Prove the chi-squared tail bound used in
Theorem~\ref{thm:highdim_jl_vector} by applying the exponential Markov
inequality to $\exp(\lambda \|g\|^2)$ and optimizing over $\lambda$. Derive
the Gaussian annulus of Theorem~\ref{thm:highdim_gaussian_annulus} from it.
\item Give an $\varepsilon$-net proof of the subspace embedding theorem
(Theorem~\ref{thm:highdim_subspace_embedding}): cover the unit sphere of a
$D$-dimensional subspace by $O((3/\varepsilon)^D)$ points, apply the
pointwise distortion bound, and take a union bound. Explain where the
exponential dependence on $D$ goes.
\item Derive the collision probability $p(u)$ of the $p$-stable hash family
of Section~\ref{sec:highdim_random_projections}, and show that it is
decreasing in $u$. For the Cauchy ($1$-stable) case compute $p(u)$ explicitly
and bound $\rho$ for the $(r, cr)$-near-neighbor problem.
\item Argue from Theorem~\ref{thm:highdim_distance_concentration} that a
$k$-nearest-neighbor classifier has error approaching the majority class on
high-dimensional Gaussian data, and design a random-projection
preprocessing that restores meaningful nearest neighbors on data lying on a
$k$-dimensional affine subspace of $\mathbb{R}^d$. Justify the target
dimension with Theorem~\ref{thm:highdim_jl_bound} and analyze the resulting
classification accuracy.
\end{enumerate}
